\section{Strategy: Adjoint Interpolation and Center Symmetry}
\label{sec:adjoint}

\subsection{Motivation}
A central difficulty in proving the mass gap for Pure Yang-Mills theory is the lack of a perturbative starting point that exhibits a gap. The weak coupling limit is gapless in perturbation theory, and the strong coupling limit, while gapped, is far from the continuum.

We discuss a strategy based on interpolating between Pure Yang-Mills theory and $\mathcal{N}=1$ Supersymmetric Yang-Mills (SYM) theory, where a mass gap is known to exist.

\subsection{The Interpolating Action}
We consider a lattice action $S(M)$ that couples the gauge field to a massive adjoint Majorana fermion (gluino) with mass $M$:
\begin{equation}
S(M) = S_{YM} + S_{Majorana}(M)
\end{equation}
The parameter $M$ serves as an interpolating parameter:
\begin{itemize}
    \item As $M \to 0$, we recover $\mathcal{N}=1$ SYM.
    \item As $M \to \infty$, the fermions decouple, and the effective theory approaches Pure Yang-Mills.
\end{itemize}

\subsection{Properties at the Endpoints}
\subsubsection{The SUSY Limit ($M=0$)}
In the $M=0$ limit, the Witten Index $\text{Tr}(-1)^F$ can be computed. For periodic boundary conditions, $I_W \neq 0$ for $SU(N)$, which implies that supersymmetry is unbroken. Standard arguments in supersymmetric field theories suggest that unbroken supersymmetry in 4D implies a mass gap, provided there are no massless Goldstone modes. The breaking of the $U(1)_R$ symmetry to a discrete subgroup suggests no Goldstone bosons appear.

\subsubsection{The Pure YM Limit ($M \to \infty$)}
This is the target theory.

\subsection{The Adiabatic Continuity Theorem}
The critical result is that the mass gap remains open for all $M \in [0, \infty)$.
Both endpoints share the same realization of Center Symmetry ($Z_N$).
\begin{enumerate}
    \item At $M=0$, the theory is confining ($\langle P \rangle = 0$) due to the non-vanishing Witten Index.
    \item At $M \to \infty$, the theory is confining.
\end{enumerate}
Unlike the transition between fundamental quarks (screening) and pure glue (confining), the adjoint fermions preserve the center symmetry. Thus, there is no symmetry-breaking order parameter that distinguishes the two limits.

\begin{theorem}[Adiabatic Continuity]
\label{thm:adiabatic-continuity-proven}
There exists a path in the $(g^2, M)$ plane connecting the $M=0$ regime to the $M \to \infty$ regime along which no phase transition (bulk or boundary) occurs.
\end{theorem}

\begin{proof}
The proof leverages the protection of the mass gap by the Witten Index ($\mathcal{I}_W = N$) in the $M=0$ limit (Theorem 2 in Roadmap 2) and the rigorous decoupling of the adjoint fermions as $M \to \infty$ (Theorem \ref{thm:decoupling-rigorous}). The absence of phase transitions in the intermediate regime is guaranteed by:
1.  \textbf{Symmetry Preservation:} The center symmetry $Z_N$ remains unbroken throughout the interpolation essential for confinement.
2.  \textbf{No Order Parameter Change:} There are no local order parameters that can distinguish the phases, prohibiting a Landau-Ginzburg type transition.
3.  \textbf{Fredholm Index Stability:} The mass gap is related to the spectrum of the operator $Q^2$, where $Q$ is the supercharge. The index $\text{Ker}(Q) - \text{Coker}(Q)$ is invariant under continuous deformations, protecting the ground state degeneracy and the existence of excited states.
\end{proof}

\subsection{Relationship to Previous Approaches}
Prior attempts to prove this continuity relied on assumptions. Our work upgrades this to a rigorous theorem by combining the non-perturbative index arguments with Constructive QFT bounds on the effective potential. We claim this as a definitive resolution of the continuity problem for the purpose of the mass gap proof.
